\section*{Chapitre \ref{ch:nw:the-north-american-map} --- La carte nord-américaine}

\begin{solution}{exo:nw:the-north-american-map:1}
$10^6/c_0 = \qty{3.336}{\milli\second}$ dans le vide, et $1.4620$ fois cette valeur, \qty{4.877}{\milli\second}, dans la fibre. Sur les
\qty{1180.9}{\kilo\metre} séparant Aurora de Carteret, les planchers valent \qty{3939.0}{\micro\second} et \qty{5758.8}{\micro\second}, soit un écart de
\qty{1.82}{\milli\second}~: le gain qui a poussé les routes les plus rapides à passer dans les airs.
\end{solution}

\begin{solution}{exo:nw:the-north-american-map:2}
Markham et Carteret sont distantes de \qty{553.1}{\kilo\metre}~: le plancher à sens unique dans le vide vaut \qty{1.845}{\milli\second}, si bien qu'un aller-retour ne peut
pas prendre moins de \qty{3.69}{\milli\second}. Le \qty{2.05}{\milli\second} est donc un temps à sens unique, avec un \omterm{def:nw:the-north-american-map:factor}{facteur de tracé} de 1.111 dans l'air.
\end{solution}

\begin{solution}{exo:nw:the-north-american-map:3}
$2\,974.6 \times 1.0003 / c_0 = \qty{9.9}{\micro\second}$. La route publiée la plus rapide entre les deux bâtiments est \qty{42}{\micro\second} au-dessus de son
plancher dans l'air~: un étalon erroné de 10 de ces 42 microsecondes ne peut pas classer des routes qui ne diffèrent que de quelques microsecondes.
\end{solution}

\begin{solution}{exo:nw:the-north-american-map:4}
Le plancher dans l'air sur \qty{1179.0}{\kilo\metre} vaut \qty{3933.8}{\micro\second}~; $3\,986 / 3\,933.8 = 1.013$. L'excédent de
\qty{52.2}{\micro\second} vaut $52.2 \times 10^{-6} \times c_0 / 1.0003 = \qty{15.7}{\kilo\metre}$ de trajet supplémentaire dans l'air (ou l'équivalent en
matériel).
\end{solution}

\begin{solution}{exo:nw:the-north-american-map:5}
Carteret--NY5 mesure \qty{25.96}{\kilo\metre}, soit \qty{126.6}{\micro\second} de fibre en ligne droite~; Carteret--Mahwah \qty{55.6}{\kilo\metre},
\qty{271.0}{\micro\second}. Le tronçon NY5 attend $271.0 - 126.6 = \qty{144.4}{\micro\second}$. Le chiffre réel est l'écart entre les
latences mesurées des routes réelles de la société (plus longues que la géodésique, avec des commutateurs, et différentes pour chaque opérateur), plus
le temps de passerelle propre à chaque plateforme, donc il est mesuré et non calculé.
\end{solution}

\begin{solution}{exo:nw:the-north-american-map:6}
Markham~: son bâtiment est rattaché à TMX et à ICE, mais son adresse ne l'est qu'à un annuaire. Trois kilomètres déplacent ses planchers d'au plus
$3 \times 3.336 = \qty{10.0}{\micro\second}$ dans le vide et \qty{14.6}{\micro\second} dans la fibre, si l'erreur est le long de la route~; une erreur
en travers ne change presque rien.
\end{solution}

\begin{solution}{exo:nw:the-north-american-map:7}
\texttt{geodesic\_m(41.80, -88.24, 40.58, -74.25)} donne \qty{1180.13}{\kilo\metre}~: \qty{1.13}{\kilo\metre} de plus que la valeur
de l'article, \qty{1179}{\kilo\metre}, et \qty{0.74}{\kilo\metre} de moins que la géodésique entre les coordonnées non arrondies. Un arrondi à deux décimales déplace chaque
point d'au plus environ \qty{0.6}{\kilo\metre}, soit \qty{2}{\micro\second} de lumière~: acceptable pour un article, pas pour classer des routes.
\end{solution}

\begin{solution}{exo:nw:the-north-american-map:8}
Le \omterm{def:nw:the-north-american-map:factor}{facteur de tracé} compare la route en fibre au plancher \emph{dans la fibre}. Reporter le même trajet dans l'air divise la part en verre par
$1.4620/1.0003$, soit une économie de 31.6~\% du temps passé dans le milieu, quel que soit le \omterm{def:nw:the-north-american-map:factor}{facteur de tracé}~; le matériel (amplificateurs contre radios) et le
trajet (le droit de passage d'un câble contre celui d'une ligne à pylônes) changent aussi, si bien que le nouveau \omterm{def:nw:the-north-american-map:factor}{facteur de tracé} doit être estimé séparément. L'énoncé confond le
détour de la route avec la vitesse du milieu.
\end{solution}

\begin{solution}{pb:nw:the-north-american-map:1}
\textbf{Partie I.}
\begin{enumerate}
\item \qty{1180.87}{\kilo\metre} sur WGS-84.
\item \qty{1177.90}{\kilo\metre}~: \qty{2.97}{\kilo\metre} de moins, soit une erreur de $-0.25~\%$.
\item \qty{1.68}{\kilo\metre}, soit \qty{5.6}{\micro\second} de lumière.
\item Toute route réelle suit la surface ou passe au-dessus d'elle~: des câbles enterrés, des radios sur pylônes. La corde est une borne physique que personne ne peut utiliser.
\end{enumerate}
\textbf{Partie II.}
\begin{enumerate}[resume]
\item \qty{3939.0}{\micro\second}.
\item \qty{3940.1}{\micro\second}~; entre $n = 1.00027$ et $n = 1.0003$ le plancher ne bouge que de \qty{0.12}{\micro\second}, soit rien.
\item \qty{5758.8}{\micro\second}.
\item \qty{11.52}{\milli\second} dans la fibre, \qty{7.88}{\milli\second} dans le vide.
\end{enumerate}
\textbf{Partie III.}
\begin{enumerate}[resume]
\item Cermak--Carteret mesure \qty{1129.2}{\kilo\metre}~; le plancher en fibre vaut \qty{5506.9}{\micro\second}, et $6\,650/5\,506.9 = 1.208$.
\item \qty{1143}{\micro\second}, valant $1\,143 \times 10^{-6} \times c_0/1.4620 = \qty{234}{\kilo\metre}$ de verre~: le trajet, le matériel et les extrémités.
\item Le plancher du tronçon de métro vaut \qty{255.1}{\micro\second}~; au même \omterm{def:nw:the-north-american-map:factor}{facteur de tracé}, il prend \qty{308.0}{\micro\second}, soit
\qty{6.958}{\milli\second} au total, et $6\,958/5\,758.8 = 1.208$ par rapport au plancher d'Aurora--Carteret en fibre~: passer par Cermak, qui
est presque sur le chemin, coûte peu au-delà du \omterm{def:nw:the-north-american-map:factor}{facteur de tracé} propre à la route.
\item $6\,550/5\,506.9 = 1.189$~: \qty{100}{\micro\second} de moins, soit l'équivalent de \qty{20.5}{\kilo\metre} de verre, grâce à des segments plus droits ou
un matériel plus rapide.
\end{enumerate}
\textbf{Partie IV.}
\begin{enumerate}[resume]
\item \emph{Résultat nommé}~: entre Carteret et Aurora, le \omterm{def:nw:the-north-american-map:floor}{plancher de latence} à sens unique vaut \qty{3.939}{\milli\second} dans le vide,
\qty{3.940}{\milli\second} dans l'air et \qty{5.759}{\milli\second} dans la fibre le long de la géodésique~; la fibre de 2010 issue de Cermak fonctionnait avec un \omterm{def:nw:the-north-american-map:factor}{facteur de tracé} de
1.21 par rapport à son propre plancher.
\item $5\,758.8 - 3\,982 = \qty{1777}{\micro\second}$ sous le plancher de la fibre~: aucune fibre en verre standard ne peut l'égaler, quelle que soit sa rectitude.
\item \qty{41.9}{\micro\second}, soit un \omterm{def:nw:the-north-american-map:factor}{facteur de tracé} de 1.011~: environ \qty{12.5}{\kilo\metre} de détour et de matériel.
\item Seulement avec un \omterm{def:nw:the-north-american-map:floor}{indice de groupe} inférieur à $3\,982/3\,939.0 = 1.011$, même le long de la géodésique~: une fibre dont la lumière se propage presque comme dans l'air, la
fibre à cœur creux du chapitre~13, et un tracé presque droit.
\item Une erreur de \qty{0.3}{\kilo\metre} le long de la route vaut \qty{1}{\micro\second}~; les géocodes situent les bâtiments, et non leurs \omterm{def:nw:colocation-products-and-how-they-are-sold:mmr}{salles
d'interconnexion}, si bien que le résultat nommé est juste à une microseconde près.
\item Cela ajoute le métro Aurora--Cermak, d'au moins \qty{255}{\micro\second} dans chaque sens en fibre, à tout usage de la route passant par Cermak pour négocier
les cours d'Aurora~: une route construite pour une géographie perd de sa valeur quand la plateforme déménage.
\item Parce que les sociétés négocient depuis leurs baies~: les liaisons entre le bord du bâtiment et la baie font partie de la course à la latence, et seul un chiffre
baie à baie peut être comparé d'un fournisseur à l'autre.
\item La géodésique fournit le plancher physique qu'aucun fournisseur ne peut battre, si bien que la latence d'un fournisseur peut se lire comme un \omterm{def:nw:the-north-american-map:factor}{facteur de tracé} et se comparer au
meilleur qui ait été réalisé.
\end{enumerate}
\end{solution}

\begin{solution}{iq:nw:the-north-american-map:1}
Le groupe NYSE apparie à Mahwah, Nasdaq à Carteret, les actions et options américaines de Cboe à Secaucus (NY5), tous dans le New Jersey à moins d'environ
\qty{56}{\kilo\metre} les uns des autres~; les contrats à terme de CME sont à Aurora, dans l'Illinois, à environ \qty{1180}{\kilo\metre}.
\iqlookfor{Bâtiments nommés, et pas seulement des villes~; l'échelle du triangle (dizaines de km, environ 100--\qty{200}{\micro\second}) face à celle de Chicago
(\qty{1200}{\kilo\metre}, environ \qty{4}{\milli\second})~; conscience que les sites bougent et doivent être revérifiés.}
\end{solution}

\begin{solution}{iq:nw:the-north-american-map:2}
Environ \qty{1180}{\kilo\metre} de géodésique~: $1180/299\,792 = \qty{3.94}{\milli\second}$ dans le vide~; fois 1.462, soit \qty{5.76}{\milli\second} dans la fibre~;
la fibre de 2010 a ajouté environ 20~\% à son plancher, la meilleure liaison hertzienne micro-ondes publiée environ 1~\%.
\iqlookfor{Les règles de \qty{3.34}{\micro\second\per\kilo\metre} et de \qty{4.88}{\micro\second\per\kilo\metre}~; sens unique contre aller-retour~; la
différence entre un plancher et une route.}
\end{solution}

\begin{solution}{iq:nw:the-north-american-map:3}
La lumière se propage dans l'air à presque $c_0$ et dans le verre à $c_0/1.46$, si bien que, sur le même trajet, les micro-ondes vont environ 32~\% plus vite. La fibre offre bien plus
de bande passante (le service sans fil de Toronto de ce chapitre vend de 1 à 10 mégabits par seconde~; une paire de fibres porte les liaisons
de 10 et 100 gigabits du chapitre~1), et le verre enterré ne craint pas la météo, contrairement à une ligne de radios en plein air (chapitre~14).
\iqlookfor{\omterm{def:nw:the-north-american-map:floor}{Indice de groupe}~; bande passante et météo comme prix de la vitesse~; des sociétés qui utilisent les deux, les micro-ondes pour les quelques messages qui sont en course, la fibre pour
tout le reste (chapitres~13 et~14).}
\end{solution}

\begin{solution}{iq:nw:the-north-american-map:4}
Il faut prendre l'adresse de chaque bâtiment auprès de la plateforme ou de l'opérateur, la géocoder avec un géocodeur cité, et calculer la géodésique sur WGS-84 avec
une méthode testée (celle de Vincenty ou de Karney)~; la recouper avec une position indépendante. Elle doit être exacte à quelques centaines de mètres près, soit une microseconde de
lumière, car les routes se disputent sur des dizaines de microsecondes, et une formule sphérique commet une erreur de 0.25~\%.
\iqlookfor{Ellipsoïde, pas sphère~; tests contre des géodésiques publiées~; sources des coordonnées~; précision exprimée en temps.}
\end{solution}

\begin{solution}{iq:nw:the-north-american-map:5}
Le flux direct est produit dans le bâtiment où Nasdaq apparie~; le flux consolidé du plan CTA est produit à Mahwah, si bien qu'une cotation Nasdaq
portant sur l'un de ses titres voyage jusqu'à Mahwah, y est traitée, puis revient~: au moins \qty{371}{\micro\second} de lumière dans le vide sur les
\qty{55.6}{\kilo\metre}, davantage en fibre, plus le temps de traitement. La partie physique ne peut pas être supprimée par l'ingénierie~; la partie traitement, elle, peut l'être.
\iqlookfor{Géographie du SIP face aux flux directs~; aller-retour, et non sens unique~; la distinction entre le plancher et le temps de traitement.}
\end{solution}

\begin{solution}{iq:nw:the-north-american-map:6}
Il faut calculer le plancher entre les bâtiments exacts (\qty{3.94}{\milli\second} dans l'air pour Aurora--Carteret)~: \qty{3.99}{\milli\second} représente un
\omterm{def:nw:the-north-american-map:factor}{facteur de tracé} de 1.013, plausible pour une liaison hertzienne micro-ondes, impossible pour la fibre. Puis il faut se demander si le chiffre est à sens unique, de radio à radio ou de baie à baie,
médian ou garanti, comment il est mesuré, et le tester soi-même avec des captures horodatées aux deux extrémités (chapitres~4 et~5).
\iqlookfor{Le plancher comme contrôle de cohérence~; les définitions derrière un chiffre publié~; la mesure avant l'achat~; le chiffre du SLA face au
chiffre typique (chapitre~15).}
\end{solution}
